\lesson{2}{Dec 06 2021 Mon (08:22:28)}{Judicial Review}{Unit 4}

\subsubsection*{Judicial Review}

Since $1789$, the \bf{Senate} has not confirmed about $\frac{1}{4}$ of \bf{Presidential Nominations} to the \bf{Supreme Court}.
In $2005$, \it{Harriet Miers}, a lawyer working for \it{President Bush}, requested that the \bf{president} withdraw his nomination of her to the
\bf{Supreme Court} amid \bf{Senate} criticism, including that \it{Miers} had never served as a judge.

\begin{definition}[Judicial Review]
    The reason why \bf{Federal Judges} are so important is because they have the power for \bf{Judicial Review}, which is the power
    of \bf{Federal Courts} to review laws of \bf{Congress} and acts of the \bf{Executive Branch} in light of the \bf{Constitution}.
    By making such a judgment, the court is declaring the law or act to be void and unenforceable—citizens are not bound to and
    \bf{Executive Agencies} may not enforce the law or act.
    
    This is a powerful check of the \bf{Federal Judiciary} on the other branches of government, since this power
    can undo \bf{Public Policy}. The \bf{Supreme Court} can check the \bf{State} and \bf{Local Governments}
    through this power as well. While the main job of the \bf{Judicial Branch} is to interpret law, the \bf{Constitution}
    does not mention \bf{Judicial Review} directly.
\end{definition}

\subsubsection*{Marbury v. Madison}

\begin{figure}[h]
    \begin{quote}
        \it{"The courts were designed to be an intermediate body between the people and the legislature,
            in order, among other things, to keep the latter within the limits assigned to their authority."}
    \end{quote}
    \caption{Alexander Hamilton, Federalist No. 78}
\end{figure}

\it{Alexander Hamilton} described the role and functions of the \bf{Judicial Branch}, including the power of the \bf{Judicial Review}.
In this quote, he names the courts as protectors of the people and a check on the \bf{Legislative Branch} because they can overrule laws
that are outside the scope of their powers under the \bf{Constitution}. While newspaper essays published during the \bf{Constitution}
ratification debates do not have the force of law, the words of \it{Hamilton} and others had a powerful effect on increasing the understanding
of the intentions of the framers.

\newpage

\begin{figure}[h]
    \begin{quote}
        \it{"The interpretation of the laws is the proper and peculiar province of the courts. A constitution is, in fact,
            and must be regarded by the judges, as a fundamental law. Where the will of the legislature, declared in its statutes,
            stands in opposition to that of the people, declared in the Constitution, the judges ought to be governed by the latter
            rather than the former. They ought to regulate their decisions by the fundamental laws, rather than by those which
            are not fundamental."}
    \end{quote}
    \caption{Alexander Hamilton, Federalist No. 78}
\end{figure}

In this quote, \it{Alexander Hamilton} explains that the \bf{Constitution} is the \it{"Fundamental Law"} of the \bf{United States},
representing the will of the people, and therefore superior in authority to any law passed by a legislative body that conflicts
with it. He explains that when judges face a conflict between a \bf{Law} and the \bf{Constitution}, they must base decisions on the Constitution.

\paragraph{Marbury v. Madison} In $1803$, the \bf{Supreme Court} decided the case of \bf{\it{Marbury VS. Madison}}. Here's the background to the case:

\begin{itemize}
    \item President John Adams appointed William Marbury to a federal position as he prepared to leave office.
    \item When President Thomas Jefferson took over, he ordered his Secretary of State James Madison not to deliver the
          commission required for Marbury to take office. Marbury was from an opposing political party.
    \item Marbury then filed a lawsuit for the court to force Madison to deliver the commission.
\end{itemize}

\it{Chief Justice John Marshall} delivered the \bf{Court Opinion} that neither the \bf{President} nor his \bf{Staff} could withhold
the commission, that to do so would violate \bf{Federal Law}.

This case may seem insignificant today, but it set an important precedent. In \bf{\it{Marbury VS. Madison}}, the justices
determined that the Supreme Court could decide whether an act or law is a violation of the Constitution. The decision
also made clear that one role of the Supreme Court is to clarify the meaning of laws and have the final say on interpretations
of the law and Constitution.

While most scholars agree that this case established the power of judicial review, they do not agree on the scope of that power.
The \bf{Constitution} is not specific on the powers of the Supreme Court. So there are questions such as whether the Supreme Court
can declare any law, statute, or act as unconstitutional or whether it is limited to certain areas. Chief Justice Marshall had a
strong role in shaping its role and powers through decisions such as Marbury v. Madison.

\begin{figure}[h]
    \begin{quote}
        \it{"It is emphatically the province and duty of the judicial department to say what the law is. A law repugnant to the Constitution is void."}
    \end{quote}
    \caption{Chief Justice John Marshall, Marbury VS. Madison ($1803$)}
\end{figure}

The first part of this quote echoes Hamilton from Federalist No. 78. Chief Justice John Marshall and his associate justices ruled in Marbury v. Madison that Madison's refusal to deliver was illegal. However, the Supreme Court did not compel the appointment nor did it have the power to do so.

The Court ruled that a portion of the law allowing Marbury to petition the Supreme Court, the Judiciary Act of 1789, was unconstitutional because it expanded jurisdiction of the court beyond areas described in Article III of the Constitution. Therefore, while the Supreme Court ruled it illegal for Madison to withhold the appointments, it could not force him to deliver them because according to the Constitution, the matter should never have come to the U.S. Supreme Court.

\subsubsection*{Judges Make Policy}

\newpage
